{\bfseries Caminata cuántica continua en un ciclo de 3 vértices. 
Sea un grafo ciclo con tres vértices etiquetados 0, 1, 2, donde cada vértice está conectado con los otros dos.  
La matriz de adyacencia es: 
\begin{equation*}
A = \begin{pmatrix} 0 & 1 & 1 \\ 1 & 0 & 1 \\ 1 & 1 & 0 \end{pmatrix} 
\end{equation*}

Considere una caminata cuántica continua con Hamiltoniano $H = A$ y estado inicial $\ket{\psi(0)} =
\ket{0}$.} 

\begin{enumerate}
    \item[a)] {\bfseries Diagonalice la matriz $H$: encuentre sus autovalores y autovectores
        ortonormales. }
    \item[b)] {\bfseries Exprese el estado inicial en la base de autovectores y escriba
        $\ket{\psi(t)}$ en términos de los autovalores. }
    \item[c)] {\bfseries Calcule las amplitudes $\bra{i}\ket{\psi(t)}$ para $i=0, 1, 2$ y las
        probabilidades $P_i(t) = |\bra{i}\ket{\psi(t)}|^2$. }
    \item[d)] {\bfseries ¿Cuál es el período de oscilación del sistema? ¿Existe algún tiempo en el
        que la partícula se localice completamente en un vértice distinto del inicial? }
\end{enumerate}
